\documentclass[11pt]{article}
\usepackage[T1]{fontenc}
\usepackage{lmodern}
\usepackage[margin=1in]{geometry}
\usepackage{amsmath,amssymb,amsthm,mathtools}
\usepackage[hidelinks]{hyperref}
\usepackage{graphicx}
\usepackage{enumitem}
\emergencystretch=1.5em

\newtheorem{theorem}{Theorem}
\newtheorem{lemma}[theorem]{Lemma}
\newtheorem{remark}[theorem]{Remark}
\newcommand{\E}{\mathbb E}
\newcommand{\R}{\mathbb R}
\newcommand{\icx}{\mathrm{icx}}

\title{A Convex-Order Solution of the\\Generalized-Cone Conjecture}
\author{Run \texttt{fa\_banach\_001}, lane 14}
\date{August 9, 2026}

\begin{document}
\maketitle

\begin{abstract}
González Merino reduced a family of sharp Hermite--Hadamard inequalities
to a one-parameter maximization over truncated cones and conjectured that
the maximum is always attained by a cone.  We prove the conjectured
inequality in every dimension.  The axial volume law of a truncated cone,
normalized by its median, has an explicit quantile.  Its quantile crosses
the cone quantile exactly once, and its mean is no larger.  These two facts
give increasing-convex-order domination by the cone law, simultaneously
for every admissible convex function.  The mean estimate is the only
nonlocal step; it follows from concavity on one side of the parameter
space and a positive-coefficient power-series ratio on the other.  For
strictly convex functions the maximizer is unique in dimensions at least
two, yielding the proposed equality case under the intended nonzero
normalization.
\end{abstract}

\section{The conjecture and the result}

Let $C\subset\R^n$ be a convex body and let
$f:C\to[0,\infty)$ be concave.  The source associates with $(C,f)$ a point
$x_{C,f}$ whose definition is based on a volume-halving section.  Its
Conjecture~1.7, reproduced in Figure~\ref{fig:question}, asks whether, for
every convex $\phi:[0,\infty)\to[0,\infty)$ with $\phi(0)=0$,
\begin{equation}\label{eq:target}
 \frac1{|C|}\int_C\phi(f(x))\,dx
 \le
 n\int_0^1
 \phi\left(f(x_{C,f})\frac{2^{1/n}t}{2^{1/n}-1}\right)
 (1-t)^{n-1}\,dt.
\end{equation}

\begin{figure}[ht]
\centering
\includegraphics[width=.94\textwidth]{figures/open_problem_crop.png}
\caption{Conjecture 1.7 in arXiv:2005.13839v2, PDF p.~5.}
\label{fig:question}
\end{figure}

\begin{theorem}[Full resolution of the inequality]\label{thm:main}
Inequality \eqref{eq:target} holds in every dimension $n\ge1$.
If $n\ge2$, $\phi$ is strictly convex, and
$f(x_{C,f})>0$, equality holds if and only if $C$ is a generalized cone
and $f$ is affine and vanishes on its nondegenerate base.
\end{theorem}

The positivity condition in the equality statement is necessary.  As the
source conjecture is literally phrased, $f\equiv0$ gives equality for every
convex body $C$, including bodies that are not generalized cones.  The
paper's optimization problem is introduced with the normalization
$f(x_{C,f})=k>0$; Theorem~\ref{thm:main} proves exactly that intended
nontrivial equality statement.

\section{Proof intuition}

The source has already done the geometric work: its main theorem replaces
an arbitrary pair $(C,f)$ by an affine function on a truncated cone.  What
remains is to compare a one-parameter family of axial probability laws.
Normalize each affine function at the volume median.  The resulting law has
an explicit quantile $R_a$.

Two elementary-looking comparisons finish the problem.  First, $R_a$ and
the cone quantile $R_0$ cross exactly once, at their common median $1$.
This follows by differentiating with respect to the truncation parameter;
the derivative sign is controlled by the decreasing elasticity of
$(1+z)^{1/n}-1$.  Second, $R_0$ has the largest mean.  On one half of the
parameter range this is the midpoint inequality for the concave function
$x^{1/n}$.  On the other half it becomes monotonicity of a ratio of two
positive power series.

The single crossing controls stop-loss functions above the median.  Below
the median, rewrite a stop-loss as the mean plus a put; the mean comparison
and the lower quantile crossing control the two terms.  Thus the cone law
dominates in increasing convex order, which is precisely the simultaneous
inequality for every admissible $\phi$.

\section{Reduction to a normalized axial law}

We use the source's Theorem~1.1 \cite[Theorem~1.1]{GM}.  It states that the left side of
\eqref{eq:target} is no larger than the maximum obtained on generalized
truncated cones $R$, with an affine function $g$ which vanishes on one
base and satisfies
\[
 g(x_{R,g})=f(x_{C,f}).
\]
Moreover, for strictly convex $\phi$, equality in this reduction forces
$C$ itself to be one of the maximizing generalized truncated cones and
$f$ to be affine and zero on the corresponding base.  It therefore
suffices to solve the remaining one-parameter maximization.

Normalize the axis of a truncated cone to $[0,1]$, put the zero-value base
at $t=0$, and scale its two base radii to $1$ and $q\ge0$.  The case in
which the first radius is zero is included as $q=\infty$.  Under normalized
volume measure, the axial coordinate $T_q$ has density proportional to
\[
 (1+(q-1)t)^{n-1},\qquad 0\le t\le1.
\]
Put
\[
 r=\frac1n,\qquad a=q^n.
\]
For $q\ne1$, the distribution function and quantile of $T_q$ are
\begin{align}
 F_q(t)
 &=\frac{(1+(q-1)t)^n-1}{q^n-1}, \label{eq:cdf}\\
 Q_a(u)
 &=\frac{(1+(a-1)u)^r-1}{a^r-1}. \label{eq:quantileT}
\end{align}
The formulas at removable values are understood by continuity.

The point $x_{R,g}$ is on the volume-halving section, so its axial
coordinate is the median $Q_a(1/2)$.  If
$c=f(x_{C,f})$, then $g(t)=ct/Q_a(1/2)$.  Thus
\[
 \frac1{|R|}\int_R\phi(g(x))\,dx=\E\phi(cX_a),
 \qquad
 X_a:=\frac{T_q}{Q_a(1/2)}.
\]
By \eqref{eq:quantileT}, the quantile of $X_a$ is
\begin{equation}\label{eq:Ra}
 R_a(u)=
 \frac{(1+(a-1)u)^r-1}{((1+a)/2)^r-1}.
\end{equation}
In particular, the cone with nondegenerate zero-value base corresponds to
$a=0$, and
\begin{equation}\label{eq:R0}
 R_0(u)=\frac{1-(1-u)^r}{1-2^{-r}}.
\end{equation}
The problem is now purely one-dimensional: prove
$X_a\le_{\icx}X_0$ for every $a\in[0,\infty]$.

\section{The quantiles cross once}

\begin{lemma}[Single crossing]\label{lem:cross}
For every $a\ge0$,
\[
 R_0(u)\le R_a(u)\quad(0\le u\le\tfrac12),
 \qquad
 R_0(u)\ge R_a(u)\quad(\tfrac12\le u\le1).
\]
\end{lemma}

\begin{proof}
For $0<r\le1$, set
\[
 \psi(z)=(1+z)^r-1,\qquad
 E(z)=\frac{z\psi'(z)}{\psi(z)}
 =\frac{rz(1+z)^{r-1}}{(1+z)^r-1}.
\]
The removable value at $z=0$ is understood continuously.  With $y=1+z$,
direct differentiation gives
\begin{equation}\label{eq:Eprime}
 E'(z)=
 \frac{r y^{r-2}}{(y^r-1)^2}
 \bigl(y^r-r y-(1-r)\bigr)\le0.
\end{equation}
Indeed, concavity of $y^r$ gives
$y^r\le r y+1-r$ for every $y\ge0$.

Set $v=(a-1)/2$ and $\tau=2u$.  Equation \eqref{eq:Ra} becomes
\[
 R_a(u)=\frac{\psi(\tau v)}{\psi(v)}.
\]
Away from removable endpoints,
\begin{equation}\label{eq:paramder}
 \frac{\partial}{\partial v}\log R_a(u)
 =\frac{E(\tau v)-E(v)}{v}.
\end{equation}
If $\tau<1$, the right side is nonnegative both for $v<0$ and $v>0$,
because $E$ is decreasing.  If $\tau>1$, it is nonpositive.  Hence
$R_a(u)$ increases with $a$ below $u=1/2$ and decreases with $a$ above
$u=1/2$.  Comparison with the endpoint $a=0$ proves the result; continuity
handles the removable cases.
\end{proof}

\section{The cone also has the largest mean}

\begin{lemma}[Mean domination]\label{lem:mean}
For every $a\in[0,\infty]$,
\begin{equation}\label{eq:meanineq}
 M(a):=\E X_a=\int_0^1R_a(u)\,du
 \le
 M(0)=\frac{r}{(r+1)(1-2^{-r})}.
\end{equation}
If $0<r<1$, equality occurs only at $a=0$.
\end{lemma}

\begin{proof}
Suppose first that $a\ge1$.  Integrating \eqref{eq:Ra} shows that
\[
 M(a)=
 \frac{\displaystyle
       \frac1{a-1}\int_1^a x^r\,dx-1}
      {\displaystyle ((1+a)/2)^r-1}.
\]
Concavity of $x^r$ gives
\[
 \frac1{a-1}\int_1^a x^r\,dx
 \le \left(\frac{1+a}{2}\right)^r,
\]
so $M(a)\le1$.  On the other hand, $M(0)\ge1$ because
$2^r\le1+r$ for $0<r\le1$.  Both inequalities are strict when
$0<r<1$ and $a>1$, while $M(1)=1<M(0)$.

It remains to consider $0\le a\le1$.  Put $h=1-a$ and
\[
 p(x)=1-(1-x)^r.
\]
For $0<r<1$ its binomial expansion on $[0,1]$ is
\[
 p(x)=\sum_{k\ge1}c_kx^k,\qquad c_k>0.
\]
Equation \eqref{eq:Ra} gives
\begin{equation}\label{eq:seriesratio}
 M(1-h)=
 \frac{\displaystyle\int_0^1p(hu)\,du}{p(h/2)}
 =
 \frac{\displaystyle\sum_{k\ge1}c_kh^k/(k+1)}
      {\displaystyle\sum_{k\ge1}c_kh^k/2^k}.
\end{equation}
This ratio is strictly increasing in $h$.  For completeness, put
\[
 A(h)=\sum_{k\ge1}\alpha_kh^k,\quad
 B(h)=\sum_{k\ge1}\beta_kh^k,\quad
 \frac{\alpha_k}{\beta_k}=\rho_k:=\frac{2^k}{k+1}.
\]
The sequence $\rho_k$ is strictly increasing, and pairing the $(i,j)$ and
$(j,i)$ terms gives
\[
 A'B-AB'
 =
 \sum_{j>i}(j-i)\beta_i\beta_j
 (\rho_j-\rho_i)h^{i+j-1}>0
 \qquad(h>0).
\]
Consequently $M(1-h)$ is largest at $h=1$, equivalently $a=0$.
Continuity gives the endpoint.  If $r=1$, \eqref{eq:Ra} reduces to
$R_a(u)=2u$ for every $a$, and the asserted non-strict inequality is
immediate.
\end{proof}

\section{Increasing-convex order and completion}

\begin{lemma}[Stop-loss domination]\label{lem:icx}
For every $a\ge0$ and every $s\ge0$,
\[
 \E(X_a-s)_+\le \E(X_0-s)_+.
\]
Consequently $X_a\le_{\icx}X_0$.
\end{lemma}

\begin{proof}
Every $X_a$ has median $1$.  If $s\ge1$, quantiles below $1/2$ contribute
nothing to either stop-loss expectation, while Lemma~\ref{lem:cross}
orders the upper quantiles in the required direction.

If $0\le s\le1$, quantiles above $1/2$ contribute nothing to the put
expectations, and the lower-half inequality in Lemma~\ref{lem:cross} gives
\[
 \E(s-X_a)_+\le\E(s-X_0)_+.
\]
Using $(x-s)_+=x-s+(s-x)_+$ and Lemma~\ref{lem:mean}, we obtain
\begin{align*}
 \E(X_0-s)_+-\E(X_a-s)_+
 &=
 \E X_0-\E X_a\\
 &\quad+\E(s-X_0)_+-\E(s-X_a)_+\ge0.
\end{align*}
This proves the stop-loss inequalities.

Finally, any finite convex nonnegative function on $[0,\infty)$ satisfying
$\phi(0)=0$ is increasing and has the call-function representation
\[
 \phi(x)=\alpha x+\int_{[0,\infty)}(x-s)_+\,\mu(ds),
 \qquad \alpha\ge0,\quad\mu\ge0,
\]
on every bounded interval.  Integrating the stop-loss inequalities proves
$\E\phi(cX_a)\le\E\phi(cX_0)$ for every $c\ge0$.
\end{proof}

\begin{proof}[Proof of Theorem~\ref{thm:main}]
For $a=0$, the axial coordinate has density
\[
 n(1-t)^{n-1},\qquad 0\le t\le1,
\]
and median
\[
 Q_0(1/2)=1-2^{-1/n}
 =\frac{2^{1/n}-1}{2^{1/n}}.
\]
Therefore
\[
 \E\phi(cX_0)
 =
 n\int_0^1
 \phi\left(c\frac{2^{1/n}t}{2^{1/n}-1}\right)
 (1-t)^{n-1}\,dt.
\]
Lemma~\ref{lem:icx} proves that this dominates every truncated-cone term
in the source's Theorem~1.1.  That theorem now gives \eqref{eq:target}.

Assume $n\ge2$ and $c>0$.  Lemma~\ref{lem:mean} is strict for every
$a>0$.  In fact the proof of Lemma~\ref{lem:icx} then gives a strict
stop-loss gap for every $0\le s\le1$.  After scaling by $c$, this gap is
strict for $0<s<c$.  The curvature measure of a strictly convex $\phi$
charges every nonempty interval, so
\[
 \E\phi(cX_a)<\E\phi(cX_0)\qquad(a>0).
\]
Thus the unique maximizing parameter is $a=0$: the truncated cone is a
cone and the affine function vanishes on its nondegenerate base.  The
equality statement in the source's Theorem~1.1 yields the claimed
characterization of $(C,f)$.

For $n=1$, $R_a(u)=2u$ for every $a$.  The scalar law is independent of
the truncation parameter, and one-dimensional convex bodies are intervals,
which are generalized cones under the natural zero-dimensional-base
convention.
\end{proof}

\section{Consistency and novelty notes}

\paragraph{The source's three-dimensional constant.}
Putting $\phi(t)=t$ in Conjecture~1.7 gives the sharp coefficient
\[
 \frac{2^{1/n}}{(n+1)(2^{1/n}-1)}.
\]
For $n=3$ this is one twelfth of the coefficient printed in the source's
Theorem~1.6.  The discrepancy is visible in the source proof: the integral
of $t(r+mt)^2$ is expanded without its denominator $12$.  The conjecture,
the planar formula, and the direct cone computation all agree with the
coefficient above.  This arithmetic issue is downstream of Theorem~1.1
and does not enter the reduction used here.

\paragraph{Literature search.}
The arXiv version and the final 2021 journal version both state the target
as Conjecture~1.7.  Bounded exact-title, formula-fragment, and core-keyword
searches performed on August 9, 2026 found no later proof or counterexample.
Crossref and OpenAlex both reported zero citing works for the published
article (DOI \texttt{10.1007/s10231-021-01080-y}).  This supports, but
cannot certify, novelty; specialist review remains appropriate.

\begin{thebibliography}{9}
\bibitem{GM}
B. González Merino,
\emph{Estimating the average of functions with convexity properties by
means of a new center},
Ann. Mat. Pura Appl. \textbf{200} (2021), 2285--2304;
arXiv:2005.13839v2.
\end{thebibliography}

\end{document}
